% Author parameterization: override with \renewcommand{\documentauthor}{Your Name} in your document
\providecommand{\documentauthor}{Nathanael Chwojko-Srawley}

\documentclass{article}
%  ╔══════════════════════════════════════════════════════════╗
%  ║                         Title                            ║
%  ╚══════════════════════════════════════════════════════════╝
%NOTE: I want \title,\author to be near top, and this \renewcommand{\maketitle} fails if these are above it.
\usepackage{titling}

\renewcommand{\maketitle}{
    \quad
    \vspace{1em}
  \begin{center}
      \LARGE\bfseries \thetitle \par
      \vspace{0.6em}
      \large
  \end{center}
    \vspace{1em}
}

\title{Spectral Sequences in Practice, and the\\ Atiyah--Hirzebruch Torsion Counterexample\\ to the Integral Hodge Conjecture}
\author{\documentauthor}
\date{\today}

%  ╔══════════════════════════════════════════════════════════╗
%  ║                         packages                         ║
%  ╚══════════════════════════════════════════════════════════╝

\usepackage{extarrows}         % Used to compile any UTF-8 Character (and supports Unicode)
\usepackage{stackengine}
\usepackage{colortbl}
\usepackage{scalerel}
\usepackage{amsmath, amssymb} % Used to write better mathematical expressions (ex. \mathbb{R})
\usepackage{stackrel}         % for left-right arrows, staking symbol above and bellow
\usepackage{mathrsfs}
\usepackage{bbm}                     % for mathbbm{1}
\usepackage{euscript}
\usepackage{commutative-diagrams}
\usepackage[font=itshape]{quoting}
\usepackage{mathtools}
  \usepackage{enumitem}
\usepackage{amsthm}           % Used to create theorem enviroments.
\usepackage{graphicx}         % Let's you use images in LaTeX; ex the \includegraphics command.
\usepackage{hhline}
%\usepackage{luatexja}
%\usepackage{fontspec}
\usepackage{dsfont}
\usepackage{wasysym}          % emoji's!
\usepackage{float}            % package with ``H'' for figures and tables, ex. Images, tables, etc
%\usepackage{subcaption}      % More options with sub-captions.
\usepackage{setspace}         % More flexibility on spacing in document.
\usepackage{hyperref}         % let's me put hyperlinks
%\usepackage{tikz}             % for commutative diagrams https://tex.stackexchange.com/questions/419221/how-to-do-the-pushout-with-universal-property
\usepackage{tikz-cd}          % for commutative diagrams https://tex.stackexchange.com/questions/419221/how-to-do-the-pushout-with-universal-property
%\usepackage{quiver}           % for better commutative diagrams
\usepackage{bookmark}         % creating heading shortcut bar on the left as in word
\usepackage{longtable}        % create tables that extend past one page
%\usepackage{siunitx}         % required for alignment
%\usepackage{multirow}        % for multiple rows
\usepackage{microtype}        % makes nice text. Fixes some under-/over- filled box problems.
%\usepackage{fancyvrb}        % Let's you write code (different style than listings).
\usepackage[most]{tcolorbox} 	      % Makes nice boxes for thms, cor, prop, or anything else.
%\usepackage{enumitem}        % More flexibility on enumeration (don't know much).
\usepackage{xcolor}           % Colour text.
\usepackage{wrapfig}          % To wrap text around figure
\usepackage{listings}         % Create nice colour code blocks (customized in preamble)
\usepackage{thmtools}         % Customize theorem environment (in preamble).
%\usepackage{marvosym}        % Provide ``basic'' symbols (hazard, religious, smiley face etc.)
\usepackage{array}            % \begin{table}{>{\em}c c} -> 1st column italic (bfseries for bold)
%\usepackage[english]{babel}  % If I want to blindtext in english.
\usepackage{blindtext}
\usepackage[a4paper, total={6in, 8.5in}]{geometry}
%\usepackage[symbol]{footmisc} % For daggers and asterisk for footnotes. 1: *, 2: dagger, etc.
\usepackage{fancyhdr}
\usepackage{titlesec}
%\usepackage{pst-plot}          %To make nice plots: https://tex.stackexchange.com/questions/47594/how-can-i-make-a-graph-of-a-function#47596
% \usepackage[answerdelayed]{exercise}
\usepackage{etoolbox}          % for Dynkin diagrams
\usepackage{dynkin-diagrams}   % for Dynkin diagrams

% ╔═════════════════════════════════════════════════════╗
% ║           for figures via inkspace                  ║
% ╚═════════════════════════════════════════════════════╝

% to import figures via: https://github.com/gillescastel/inkscape-figures
\usepackage{import}
\usepackage{pdfpages}
\usepackage{transparent}


% This command is the ONLY command imported via the packages snippet, think of it as a tiny package written out
\newcommand{\incfig}[2][1]{%
    \def\svgwidth{#1\columnwidth}
    \import{./figures/}{#2.pdf_tex}
}

% ╔═══════════════════════════════════════════════════════════╗
% ║                 bibliography and citing                   ║
% ╚═══════════════════════════════════════════════════════════╝

\usepackage{csquotes}
% \usepackage[backend=biber,citestyle=alphabetic,bibstyle=authortitle]{biblatex}
\usepackage{biblatex}

\addbibresource{bibliography.bib}

% ╔═══════════════════════════════════════════════════════════╗
% ║                    colours (colors)                       ║
% ╚═══════════════════════════════════════════════════════════╝

\definecolor{dkgreen}{rgb}{0,0.6,0}
\definecolor{gray}{rgb}{0.5,0.5,0.5}
\definecolor{mauve}{rgb}{0.58,0,0.82}
\definecolor{lightyellow}{rgb}{1,1,0.6}
\definecolor{reddish}{HTML}{A01A3A}

%  ╔══════════════════════════════════════════════════════════╗
%  ║                       book format                        ║
%  ╚══════════════════════════════════════════════════════════╝

\pagestyle{fancy}

\setlength\headheight{20pt}

\renewcommand{\sectionmark}[1]{\markright{\thesection.\ #1}}

\fancyfoot[C,CO]{\textcolor{gray}{\documentauthor}}
\fancyfoot[R,RO]{\thepage}{}

% Create a new page style for the first page
\fancypagestyle{firstpage}{
  \fancyhf{} % Clear header and footer
  \fancyhead[L]{\theauthor}
  \fancyhead[R]{\today} % Date in the top-right
  \fancyfoot[C]{\textcolor{gray}{\documentauthor}}
  \fancyfoot[R]{\thepage}
}

% Apply the first page style conditionally
\AtBeginDocument{\thispagestyle{firstpage}}

%have format the title command
\newcommand{\chapter}[1]{%
  \clearpage
  \vspace*{30pt} % Space before the chapter title
  {\normalfont\huge\bfseries #1\par} % Chapter title formatting
  \vspace{20pt} % Space after the chapter title
}


\setlength{\parindent}{0pt} % don't like indentation infront of paragraphs
\setlength{\parskip}{1ex}   %so that there is still space between paragarphs


%  ╔══════════════════════════════════════════════════════════╗
%  ║                       Shortcuts                          ║
%  ╚══════════════════════════════════════════════════════════╝

%\ExerciseLevelInToc

% I like original mathcal better, but need lowercase.
\DeclareFontFamily{OT1}{pzc}{}
\DeclareFontShape{OT1}{pzc}{m}{it}{<-> s * [1.10] pzcmi7t}{}
\DeclareMathAlphabet{\mathpzc}{OT1}{pzc}{m}{it}


% mathbb
\newcommand{\A}{{\mathbb{A}}}
\newcommand{\C}{{\mathbb{C}}}
\newcommand{\F}{{\mathbb{F}}}
\newcommand{\N}{{\mathbb{N}}}
\newcommand{\Q}{{\mathbb{Q}}}
\newcommand{\R}{{\mathbb{R}}}
\newcommand{\V}{{\mathbb{V}}}
\newcommand{\Z}{{\mathbb{Z}}}
\newcommand{\BI}{{\mathbb{I}}}
\renewcommand{\H}{{\mathbb{H}}}
\renewcommand{\P}{{\mathbb{P}}}

\DeclareMathOperator{\Sq}{Sq}
\DeclareMathOperator{\CH}{CH}
\DeclareMathOperator{\Hdg}{Hdg}
\newcommand{\RP}{\mathbb{RP}}
\newcommand{\CP}{\mathbb{CP}}
\newcommand{\KU}{\mathrm{KU}}
\newcommand{\ptt}{\mathrm{pt}}
\renewcommand{\O}{\mathcal{O}}
\newcommand{\gr}{\operatorname{gr}}
% \newcommand{\to}{\longrightarrow}


% mathcal
% \newcommand{\co}{{\mathfrak{o}}} %mathcal{o} looks like wreath product, so I changed it to mathfrak
\newcommand{\co}{{\mathpzc{o}}}
\newcommand{\CA}{{\mathcal{A}}}
\newcommand{\CB}{{\mathcal{B}}}
\newcommand{\CC}{{\mathcal{C}}}
\newcommand{\CD}{{\mathcal{D}}}
\newcommand{\CF}{{\mathcal{F}}}
\newcommand{\CG}{{\mathcal{G}}}
\newcommand{\CI}{{\mathcal{I}}}
\newcommand{\CK}{{\mathcal{K}}}
\newcommand{\CL}{{\mathcal{L}}}
\newcommand{\CM}{{\mathcal{M}}}
\newcommand{\CN}{{\mathcal{N}}}
\newcommand{\CO}{{\mathcal{O}}}
\newcommand{\CQ}{{\mathcal{Q}}}
\newcommand{\CR}{{\mathcal{R}}}
\newcommand{\CS}{{\mathcal{S}}}
\newcommand{\CT}{{\mathcal{T}}}
\newcommand{\CV}{{\mathcal{V}}}
\newcommand{\CZ}{{\mathcal{Z}}}
% EXCPETION:
% \newcommand{\CP}{\mathbb{C}\mathrm{P}}


%Euscript
\newcommand{\EB}{\EuScript{B}}
\newcommand{\EC}{{\EuScript{C}}}
\newcommand{\EF}{{\EuScript{F}}}
\newcommand{\EL}{{\EuScript{L}}}
\newcommand{\EO}{{\EuScript{O}}}


%mathfrak (lowercase)
\newcommand{\fa}{{\mathfrak{a}}}
\newcommand{\fb}{{\mathfrak{b}}}
\newcommand{\fc}{{\mathfrak{c}}}
\newcommand{\fd}{{\mathfrak{d}}}
\newcommand{\fm}{{\mathfrak{m}}}
\newcommand{\fn}{{\mathfrak{n}}}
\newcommand{\fo}{{\mathfrak{o}}}
\newcommand{\fp}{{\mathfrak{p}}}
\newcommand{\fq}{{\mathfrak{q}}}


%mathfrak (upper case)
\newcommand{\FA}{{\mathfrak{A}}}
\newcommand{\FB}{{\mathfrak{B}}}
\newcommand{\FC}{{\mathfrak{C}}}
\newcommand{\FD}{{\mathfrak{D}}}
\newcommand{\FK}{{\mathfrak{K}}}
\newcommand{\FM}{{\mathfrak{M}}}
\newcommand{\FN}{{\mathfrak{N}}}
\newcommand{\FO}{{\mathfrak{O}}}
\newcommand{\FP}{{\mathfrak{P}}}


%mathscr
\newcommand{\ff}{{\mathscr{F}}}
\newcommand{\SA}{\mathscr{A}}
\newcommand{\SC}{\mathscr{C}}
\newcommand{\SF}{\mathscr{F}}
\newcommand{\SG}{\mathscr{G}}
\newcommand{\SH}{\mathscr{H}}
\newcommand{\SI}{\mathscr{I}}
\newcommand{\SK}{\mathscr{K}}
\newcommand{\SN}{\mathscr{N}}
\newcommand{\SQ}{\mathscr{Q}}
\newcommand{\SR}{\mathscr{R}}
\newcommand{\SV}{\mathscr{V}}
\newcommand{\SZ}{\mathscr{Z}}


% operatorname -- 2 letters (lowercase and upper case)
\newcommand{\ab}{{\operatorname{ab}}}
\newcommand{\ad}{{\operatorname{ad}}}
\newcommand{\ch}{{\operatorname{ch}}}
\newcommand{\ev}{{\operatorname{ev}}}
\newcommand{\id}{{\operatorname{id}}}
\newcommand{\im}{{\operatorname{im}}}
\newcommand{\nm}{{\operatorname{Nm}}}
\newcommand{\op}{{\operatorname{op}}}
\newcommand{\rk}{{\operatorname{rk}}}
\newcommand{\sh}{{\operatorname{sh}}}
\newcommand{\tr}{{\operatorname{tr}}}

\newcommand{\Ab}{{\operatorname{Ab}}}
\newcommand{\Br}{{\operatorname{Br}}}
\newcommand{\Ch}{{\operatorname{Ch}}}
\newcommand{\Cl}{{\operatorname{Cl}}}
\newcommand{\GL}{{\operatorname{GL}}}
\newcommand{\Nm}{{\operatorname{Nm}}}
\newcommand{\SL}{{\operatorname{SL}}}

\renewcommand{\Re}{{\operatorname{Re}}}
\renewcommand{\Im}{{\operatorname{Im}}}


%categories
\newcommand{\Grp}{{\textbf{Grp}}}
\newcommand{\Fld}{{\textbf{Fld}}}
\newcommand{\Sh}{{\textbf{Sh}}}
\newcommand{\Set}{{\textbf{Set}}}
\newcommand{\Mod}{{\textbf{Mod}}}
\newcommand{\PSh}{{\textbf{PSh}}}
\newcommand{\Sch}{{\textbf{Sch}}}
\newcommand{\Top}{{\textbf{Top}}}
\newcommand{\RgSp}{{\textbf{RgSp}}}
\newcommand{\Ring}{{\textbf{Ring}}}


%operatorname -- 3+ letters (lowercase and upper case)
\newcommand{\rad}{{\operatorname{rad}}}
\newcommand{\sep}{{\operatorname{sep}}}
\newcommand{\res}{{\operatorname{res}}}
\newcommand{\sgn}{{\operatorname{sgn}}}
\newcommand{\pre}{{\operatorname{pre}}}
\newcommand{\ord}{{\operatorname{ord}}}
\newcommand{\vol}{{\operatorname{vol}}}
\newcommand{\lcm}{{\operatorname{lcm}}}

\newcommand{\conj}{{\operatorname{conj}}}
\newcommand{\disc}{{\operatorname{disc}}}
\newcommand{\stab}{{\operatorname{stab}}}
\newcommand{\kdim}{{\operatorname{kdim}}}
\newcommand{\diag}{{\operatorname{diag}}}

\newcommand{\trdeg}{\operatorname{trdeg}}
\newcommand{\coker}{{\operatorname{coker}}}
\newcommand{\codim}{{\operatorname{codim}}}

\newcommand{\Ann}{{\operatorname{Ann}}}
\newcommand{\Aut}{{\operatorname{Aut}}}
\newcommand{\Der}{{\operatorname{Der}}}
\newcommand{\Div}{{\operatorname{Div}}}
\newcommand{\Emb}{{\operatorname{Emb}}}
\newcommand{\End}{{\operatorname{End}}}
\newcommand{\Ext}{{\operatorname{Ext}}}
\newcommand{\Gal}{{\operatorname{Gal}}}
\newcommand{\Hom}{{\operatorname{Hom}}}
\newcommand{\Ind}{{\operatorname{Ind}}}
\newcommand{\Inn}{{\operatorname{Inn}}}
\newcommand{\Int}{{\operatorname{int}}}
\newcommand{\Irr}{{\operatorname{Irr}}}
\newcommand{\Jac}{{\operatorname{Jac}}}
\newcommand{\Mor}[1]{\operatorname{Mor}(\textbf{#1})}
\newcommand{\Nil}{{\operatorname{Nil}}}
\newcommand{\Obj}{{\operatorname{Obj}}}
\newcommand{\Out}{{\operatorname{Out}}}
\newcommand{\Pic}{{\operatorname{Pic}\,}}
\newcommand{\Rep}{{\operatorname{Rep}}}
\newcommand{\Res}{{\operatorname{Res}}}
\newcommand{\Spl}{{\operatorname{Spl}}}
\newcommand{\Syl}{{\operatorname{Syl}}}
\newcommand{\Sym}{{\operatorname{Sym}}}
\newcommand{\Tor}{{\operatorname{Tor}}}

\newcommand{\Char}{{\operatorname{char}}}
\newcommand{\Frac}{{\operatorname{Frac}}}
\newcommand{\Frob}{{\operatorname{Frob}}}
\newcommand{\Proj}{{\operatorname{Proj}\,}}
\newcommand{\RMod}{{}_R\operatorname{Mod}}
\newcommand{\SExt}{{\mathcal{E}\emph{xt}}}
\newcommand{\SHom}{{\emph{Hom}}}
\newcommand{\Span}{{\operatorname{span}}}
\newcommand{\Spec}{{\operatorname{Spec}\,}}
\newcommand{\Supp}{{\operatorname{Supp}}}
\newcommand{\Tens}[1]{#1\otimes --}
\newcommand{\fHom}[1]{\operatorname{Hom}(#1, --)}

\newcommand{\mSpec}{{\operatorname{mSpec}}}
\newcommand{\cofHom}[1]{\operatorname{Hom}(--, #1)}


% connectors
% \renewcommand{\mid}{\big|}
\newcommand{\sse}{\subseteq}
\newcommand{\subg}{\leqslant}
\newcommand{\supg}{\geqslant}
\newcommand{\ssne}{\subsetneq}
\newcommand{\isosubg}{\lesssim}
\newcommand{\nsse}{\not\subseteq}
\newcommand{\nsubg}{\trianglelefteq}
\newcommand{\nsupg}{\trianglerighteq}
\newcommand{\csubg}{\blacktriangleleft}
\renewcommand{\mid}{\big|}


% Arrows
\newcommand{\rw}{\rightarrow}
\newcommand{\Rw}{\Rightarrow}
\newcommand{\lw}{\leftarrow}
\newcommand{\Lw}{\Leftarrow}
\newcommand{\lrw}{\leftrightarrow}
\newcommand{\LRw}{\Leftrightarrow}
\newcommand{\acton}{\curvearrowright}
\newcommand{\actson}{\curvearrowright}
\newcommand\mapsfrom{\mathrel{\reflectbox{\ensuremath{\mapsto}}}}

% arrows for commutative diagram: create four new angled double-struck arrows
\newcommand\myrot[1]{\mathrel{\rotatebox[origin=c]{#1}{$\Rightarrow$}}}
\newcommand\NEarrow{\myrot{45}}
\newcommand\NWarrow{\myrot{135}}
\newcommand\SWarrow{\myrot{-135}}
\newcommand\SEarrow{\myrot{-45}}


%shorthands
\newcommand{\oln}[1]{{\overline{#1}}}
\renewcommand{\bf}[1]{\textbf{#1}}
\newcommand{\qn}{\quad\newline}


% limits
\newcommand{\Lim}{{\varprojlim}}
\newcommand{\coLim}{{\varinjlim}}
\newcommand{\Dlim}{\lim\limits_{\longrightarrow}}
\newcommand{\coDlim}{\lim\limits_{\longleftarrow}}


%for saturated ideals in the localization section (I don't like these, I think I'll delete)
\newcommand{\LIm}{{}^e }
\newcommand{\LPre}{{}^c }

%  ╔══════════════════════════════════════════════════════════╗
%  ║                    Custom Commands                       ║
%  ╚══════════════════════════════════════════════════════════╝

% swap varphi and phi
\let\temp\phi
\let\phi\varphi
\let\varphi\temp

\newcommand{\placeholder}{\_\_\_\_\_}

\newcommand{\set}[2]{\left\{#1 \ \middle|\ #2\right\}}

%mod should not have the crazy amount of space to its right!
\renewcommand{\mod}{\,\operatorname{mod}\,}

% dcups
\newcommand{\dcup}{\sqcup}
\newcommand{\Dcup}{\coprod}
% \newcommand{\Dcup}{\bigsqcup}

\newcommand{\nbot}{\mathrel{\rlap{$\bot$}\mkern2mu/}}

%a nice large circle
\newcommand{\tikzcircle}[2][red,fill=black]{\tikz[baseline=-0.5ex]\draw[#1,radius=#2] (0,0) circle ;}%

\makeatletter
\newcommand*\bigcdot{\mathpalette\bigcdot@{.5}}
\newcommand*\bigcdot@[2]{\mathbin{\vcenter{\hbox{\scalebox{#2}{$\m@th#1\bullet$}}}}}
\makeatother

%somehow not a default command
\def\rddots{\mathstrut^{.^{.^{.}}}}

%% new delimiter
\DeclarePairedDelimiter{\ceil}{\lceil}{\rceil}


%  ╔══════════════════════════════════════════════════════════╗
%  ║                     environments                         ║
%  ╚══════════════════════════════════════════════════════════╝


%remember the option: breakable

% create tcolor boxes
\newtcbtheorem[number within=section]{thm}{Theorem}%
{
colback=green!5,
colframe=green!35!black,
fonttitle=\bfseries,
label type=theorem
}{th}

\newtcbtheorem[number within=section, use counter from=thm]{lem}{Lemma}%
{
colback=blue!5,
colframe=black!35!black,
fonttitle=\bfseries,
label type=lemma
}{lm}
\newtcbtheorem[number within=section, use counter from=thm]{prop}{Proposition}%
{
colback=white!5,
colframe=white!35!black,
fonttitle=\bfseries,
label type=proposition
}{pr}
\newtcbtheorem[number within=section, use counter from=thm]{cor}{Corollary}%
{colback=blue!5,
colframe=blue!35!black,
fonttitle=\bfseries,
label type=corollary
}{co}
\newtcbtheorem[number within=section, use counter from=thm]{axiom}{Axiom}%
{colback=black!5,
colframe=yellow!35!black,
fonttitle=\bfseries,
label type=axiom
}{ax}
\newtcbtheorem[number within=section, use counter from=thm]{defn}{Definition}%
{colback=red!5,
colframe=red!35!black,
fonttitle=\bfseries,
label type=definition
}{df}


% for <s-cr> to create \item[$\LRw$] instead of \item
\newenvironment{equivEnumerate}
 {\begin{enumerate}
	}{
\end{enumerate}}



%Custom enviroments
 \newenvironment{note}
 {\begin{tcolorbox}[enhanced, sharp corners, colback=white , borderline={0.2pt}{0pt}{black}]
   \begin{quote}\textbf{Note}:
   }{
   \end{quote}\end{tcolorbox}}

 \newenvironment{intuition}
 {\begin{quote}\textbf{Intuition}:
   }{
   \end{quote}}

\newtcbtheorem[number within=section, use counter from=thm]{titledBox}{}{%
  enhanced,
  sharp corners,
  colback=white,
  colframe=black,
  borderline={0.2pt}{0pt}{black},
  left=2mm,
  right=2mm,
  top=2mm,
  bottom=2mm,
  title style={opacity=1},
  attach title to upper,
  before upper={\textbf{\tcbtitletext}\quad},
  coltitle=black,
  fonttitle=\bfseries,
  separator sign={\quad}
}{box}


%better example and proof environments
\def\exampletext{Example} % If English
\colorlet{colexam}{red!55!black} % Global example color


\newtcbtheorem[number within=section, use counter from=thm]{example}{Example}{% \newtcolorbox[use counter=testexample]{testexamplebox}{%
% Example Frame Start
empty,% Empty previously set parameters
title={\exampletext \thetcbcounter #1},% use \thetcbcounter to access the testexample counter text
% Attaching a box requires an overlay
attach boxed title to top left,
   % Ensures proper line breaking in longer titles
   minipage boxed title,
% (boxed title style requires an overlay)
boxed title style={empty,size=minimal,toprule=0pt,top=4pt,left=3mm,overlay={}},
coltitle=colexam,fonttitle=\bfseries,
before=\par\medskip\noindent,parbox=false,boxsep=0pt,left=3mm,right=0mm,top=2pt,breakable,pad at break=0mm,
   before upper=\csname @totalleftmargin\endcsname0pt, % Use instead of parbox=true. This ensures parskip is inherited by box.
% Handles box when it exists on one page only
overlay unbroken={\draw[colexam,line width=.5pt] ([xshift=-0pt]title.north west) -- ([xshift=-0pt]frame.south west); },
% Handles multipage box: first page
overlay first={\draw[colexam,line width=.5pt] ([xshift=-0pt]title.north west) -- ([xshift=-0pt]frame.south west); },
% Handles multipage box: middle page
overlay middle={\draw[colexam,line width=.5pt] ([xshift=-0pt]frame.north west) -- ([xshift=-0pt]frame.south west); },
% Handles multipage box: last page
overlay last={\draw[colexam,line width=.5pt] ([xshift=-0pt]frame.north west) -- ([xshift=-0pt]frame.south west); }%
}{ex}



\def\prooftext{Proof} % If English
\NewDocumentEnvironment{Proof}{ O{} }
{
    \colorlet{colexam}{black!55!black} % Global example color
    \newtcolorbox[number within=section]{testproofbox}{%
        empty,% Empty previously set parameters
        title={\emph{\prooftext} \thetcbcounter: #1},
        attach boxed title to top left,
        minipage boxed title,
        boxed title style={empty,size=minimal,toprule=0pt,top=4pt,left=3mm,overlay={}},
        coltitle=colexam,
        fonttitle=\bfseries,
        before=\par\medskip\noindent,
        parbox=false,
        boxsep=0pt,
        left=3mm,
        right=0mm,
        top=2pt,
        breakable,
        pad at break=0mm,
        before upper=\csname @totalleftmargin\endcsname0pt,
        % Removed height constraints
        % Added flexible width
        width=\dimexpr\textwidth-3mm\relax,
        % Modified overlays
        overlay unbroken={\draw[colexam,line width=.5pt] ([xshift=-0pt]title.north west) -- ([xshift=-0pt]frame.south west); },
        overlay first={\draw[colexam,line width=.5pt] ([xshift=-0pt]title.north west) -- ([xshift=-0pt]frame.south west); },
        overlay middle={\draw[colexam,line width=.5pt] ([xshift=-0pt]frame.north west) -- ([xshift=-0pt]frame.south west); },
        overlay last={\draw[colexam,line width=.5pt] ([xshift=-0pt]frame.north west) -- ([xshift=-0pt]frame.south west); },%
    }
    \begin{testproofbox}
}
{\end{testproofbox}\endlist}


%  ╔══════════════════════════════════════════════════════════╗
%  ║                    for Dynkin diagrams                   ║
%  ╚══════════════════════════════════════════════════════════╝

\def\row#1/#2!{#1_{\IfStrEq{#2}{}{n}{#2}} & \dynkin{#1}{#2}\\}
\newcommand{\tble}[1]{
   \renewcommand*\do[1]{\row##1!}
   \[
      \begin{array}{ll}\docsvlist{#1}\end{array}
   \]
}

%  ╔══════════════════════════════════════════════════════════╗
%  ║                         links                            ║
%  ╚══════════════════════════════════════════════════════════╝

\definecolor{LinkColour}{HTML}{A64A3B} % favourite
% \definecolor{LinkColour}{HTML}{8C3D32} % 2nd place
\hypersetup{
  colorlinks,
  citecolor=black,
  filecolor=magenta,
  linkcolor=LinkColour,
  urlcolor=blue,
}


%  ╔══════════════════════════════════════════════════════════╗
%  ║                          misc                            ║
%  ╚══════════════════════════════════════════════════════════╝

\stackMath %to be able to stack text in math mode
\makeindex

%import options: all, bibliography, book-formatting, booksSetting, customEnvironments, dynkin, hw-formatting, language-setup, newcommands, packages, preamble, proofs, settings, tcolorbox, test.svg,
\begin{document}
\maketitle
\tableofcontents

%======================================================================
\section{Conventions on spectral sequences}
%======================================================================

A (cohomological, first--quadrant) spectral sequence is a sequence of bigraded
abelian groups $\{E_r^{p,q}\}$ together with differentials
\[
  d_r \colon E_r^{p,q}\to E_r^{p+r,\,q-r+1}, \qquad d_r\circ d_r=0,
\]
such that $E_{r+1}^{p,q}=\ker d_r/\operatorname{im} d_r$ at $(p,q)$. We say it
\emph{converges} to a graded group $H^\bullet$, written $E_2^{p,q}\Rightarrow
H^{p+q}$, if $H^n$ carries a finite decreasing filtration
$H^n=F^0\supseteq F^1\supseteq\cdots\supseteq F^{n}\supseteq F^{n+1}=0$ with
\[
  E_\infty^{p,q}\cong F^p H^{p+q}/F^{p+1}H^{p+q}.
\]
Three separate pieces of information are always at play, and it is worth naming
them because \emph{each} is a distinct step in a computation:

\begin{enumerate}
\item[(D)] \textbf{Differentials.} One must determine the maps $d_r$. Often they
are forced by knowing the abutment (as in Part~I), or identified with a known
cohomology operation (as in Part~II).
\item[(M)] \textbf{Multiplicativity.} In good cases the $E_r$ are bigraded rings,
$d_r$ is a derivation
$d_r(xy)=d_r(x)\,y+(-1)^{|x|}x\,d_r(y)$ (with $|x|=p+q$), and the filtration on
$H^\bullet$ is multiplicative. This propagates a single differential across the
whole page.
\item[(E)] \textbf{The extension problem.} Even after computing $E_\infty$, one
only knows the \emph{associated graded} $\gr H^n=\bigoplus_p E_\infty^{p,n-p}$.
Recovering $H^n$ requires solving extensions
$0\to F^{p+1}\to F^{p}\to E_\infty^{p,n-p}\to 0$, which can be nonsplit.
\end{enumerate}

Parts~I and~III(the $\RP^{2n}$ example) are chosen precisely so that you meet all
three phenomena.

%======================================================================
\section{A fully worked example: the Serre SS of a path--loop fibration}
%======================================================================

\subsection{The setup}

Fix $n\ge 2$ and let $PS^n=\{\gamma\colon[0,1]\to S^n : \gamma(0)=x_0\}$ be the
space of based paths, with the fibration
\[
  \Omega S^n \hookrightarrow PS^n \xrightarrow{\ \mathrm{ev}_1\ } S^n,
  \qquad \mathrm{ev}_1(\gamma)=\gamma(1),
\]
whose fibre is the based loop space $\Omega S^n$. The total space $PS^n$ is
\emph{contractible} (contract each path to the constant path), so
\begin{equation}\label{eq:contractible}
  H^m(PS^n;\Z)=\begin{cases}\Z,& m=0,\\ 0,&m>0.\end{cases}
\end{equation}
Since $S^n$ is simply connected for $n\ge 2$, the local system on the base is
trivial and the cohomology Serre spectral sequence reads
\[
  E_2^{p,q}=H^p\!\big(S^n;\,H^q(\Omega S^n;\Z)\big)
  \;\Longrightarrow\; H^{p+q}(PS^n;\Z).
\]
Write $A^q:=H^q(\Omega S^n;\Z)$; this is what we want to compute. Because
$H^p(S^n;\Z)=\Z$ for $p\in\{0,n\}$ and $0$ otherwise, \emph{only two columns are
nonzero}:
\[
  E_2^{0,q}=A^q,\qquad E_2^{n,q}=A^q,\qquad E_2^{p,q}=0 \text{ for } p\neq 0,n.
\]

\subsection{Which differentials can be nonzero?}

A differential $d_r\colon E_r^{p,q}\to E_r^{p+r,q-r+1}$ can be nonzero only if
both columns $p$ and $p+r$ are nonzero, i.e.\ only if $\{p,p+r\}=\{0,n\}$. Hence
\emph{the only possibly nonzero differential is $d_n$}, and
\[
  E_2=E_3=\cdots=E_n,\qquad E_{n+1}=E_\infty .
\]
The single relevant family of maps is the \emph{transgression}
\[
  d_n\colon E_n^{0,q}=A^{q}\;\to\;E_n^{n,q-n+1}=A^{q-n+1}.
\]

\subsection{Convergence forces $d_n$ to be an isomorphism}

By \eqref{eq:contractible} the abutment is trivial except in total degree $0$, so
$E_\infty^{p,q}=0$ unless $(p,q)=(0,0)$. Since $E_{n+1}=E_\infty$, the only
group that may survive is $E_\infty^{0,0}=\Z$.

\smallskip\noindent\textbf{Low degrees.} For $0<q<n-1$ the group $E_n^{0,q}=A^q$
receives no differential (its potential source $E_n^{n,q-n+1}$ sits in negative
fibre degree) and maps nowhere nontrivially; hence it survives to $E_\infty$ and
must vanish:
\[
  A^q=0\qquad(1\le q\le n-2).
\]

\smallskip\noindent\textbf{Transgression.} In degree $q=n-1$ the class
$s\in H^n(S^n)=E_n^{n,0}$ must die. The only differential that can kill it is
\[
  d_n\colon E_n^{0,n-1}=A^{n-1}\xrightarrow{\ \cong\ }E_n^{n,0}=\Z\langle s\rangle .
\]
For $E_\infty^{n,0}=0$ this map must be surjective; for $E_\infty^{0,n-1}=0$ it
must be injective. Therefore $A^{n-1}=\Z$, generated by a class $a$ with
\begin{equation}\label{eq:transgress}
  d_n(a)=s .
\end{equation}

\smallskip\noindent\textbf{Bootstrapping.} More generally, for every $q$ the map
$d_n\colon A^{q}\to A^{q-n+1}$ must be an isomorphism whenever both are off the
corner (source in column $0$, target in column $n$, and neither is
$E_\infty^{0,0}$). Feeding in $A^0=\Z$, $A^{q}=0$ for $1\le q\le n-2$, and
iterating $A^{q}\cong A^{q-(n-1)}$ gives:
\begin{equation}\label{eq:additive}
  H^{q}(\Omega S^n;\Z)=
  \begin{cases}\Z,& q\equiv 0\pmod{n-1},\ q\ge 0,\\[2pt] 0,&\text{otherwise.}\end{cases}
\end{equation}
Denote the generator of $H^{k(n-1)}$ by $a_k$, with $a_0=1$ and $a_1=a$.

\subsection{A concrete page: $n=3$}

For $S^3$ the fibre cohomology lives in even degrees $0,2,4,6,\dots$, and the
$E_2=E_3$ page has just the two columns $p=0,3$:
\[
\renewcommand{\arraystretch}{1.35}
\begin{array}{c|cccc}
 q=6 & \Z\langle a_3\rangle & 0 & 0 & \Z\langle s\,a_3\rangle \\
 q=4 & \Z\langle a_2\rangle & 0 & 0 & \Z\langle s\,a_2\rangle \\
 q=2 & \Z\langle a_1\rangle & 0 & 0 & \Z\langle s\,a_1\rangle \\
 q=0 & \Z\langle 1\rangle   & 0 & 0 & \Z\langle s\rangle \\ \hline
     & p=0 & 1 & 2 & p=3
\end{array}
\]
The differential $d_3$ points three columns to the right and two rows down,
$d_3\colon E_3^{0,2k}\to E_3^{3,2k-2}$, i.e.
\[
  a_k\;\longmapsto\; (\text{integer})\cdot s\,a_{k-1}.
\]
Convergence (everything off the corner must die) forces each such $d_3$ to be an
\emph{isomorphism} $\Z\xrightarrow{\cong}\Z$; after fixing signs we may take
\begin{equation}\label{eq:d3loop}
  d_3(a_k)=s\,a_{k-1}\qquad(k\ge 1).
\end{equation}

\subsection{Multiplicativity and the integral extension problem}

So far \eqref{eq:additive} is purely additive. Now use that the Serre spectral
sequence is \emph{multiplicative}: $E_2^{\bullet,\bullet}=H^*(S^n)\otimes
H^*(\Omega S^n)$ as bigraded rings, and $d_n$ is a derivation. Take $n=3$, so
$a=a_1$ has even total degree and there are no Koszul signs. From
\eqref{eq:d3loop} and the Leibniz rule,
\[
  d_3(a_1^2)=d_3(a_1)\,a_1+a_1\,d_3(a_1)=s\,a_1+a_1\,s=2\,s\,a_1 .
\]
But $a_1^2\in H^4(\Omega S^3)=\Z\langle a_2\rangle$, say $a_1^2=\lambda a_2$, and
by \eqref{eq:d3loop} $d_3(a_2)=s\,a_1$. Hence
\[
  2\,s\,a_1=d_3(a_1^2)=\lambda\,d_3(a_2)=\lambda\,s\,a_1
  \quad\Longrightarrow\quad \boxed{\,a_1^2=2\,a_2\,}.
\]
This is the crux: the ring is \emph{not} polynomial. Iterating (induction with
$a_1a_{k-1}=k\,a_k$, obtained the same way) gives
\[
  a_1^{\,k}=k!\;a_k .
\]

\begin{thm}{Loop}{loop}
For $n$ odd, $H^*(\Omega S^n;\Z)$ is the \emph{divided power algebra}
$\Gamma_\Z[a]$ on one generator $a$ of degree $n-1$: it is free of rank one in
each degree $k(n-1)$ on a class $\gamma_k$, with $\gamma_1=a$,
$a^{k}=k!\,\gamma_k$, and $\gamma_i\gamma_j=\binom{i+j}{i}\gamma_{i+j}$. In
particular it is \emph{not} a polynomial ring over $\Z$. \qed
\end{thm}

\begin{titledBox}{Note}{noteI}
Every technique you need downstream appeared here:
\textbf{(D)} differentials forced by the abutment via \eqref{eq:contractible};
\textbf{transgression} \eqref{eq:transgress};
\textbf{(M)} a single differential propagated by the Leibniz rule;
\textbf{(E)} an extension/divided‑power subtlety visible only integrally
($a_1^2=2a_2$). Reducing mod a prime $p$ would hide this: $H^*(\Omega S^3;\F_p)$
\emph{is} polynomial. The moral --- ``torsion and integral structure hide in the
multiplicative/extension data'' --- is exactly what powers Part~IV.
\end{titledBox}

%======================================================================
\section{The Atiyah--Hirzebruch spectral sequence}
%======================================================================

\subsection{Definition}

Let $E^\bullet$ be a generalized (multiplicative) cohomology theory and $X$ a
finite CW complex, filtered by skeleta $X^{(p)}$. The exact couple of the
filtration produces the \emph{Atiyah--Hirzebruch spectral sequence} (AHSS)
\begin{equation}\label{eq:AHSS}
  E_2^{p,q}=H^p\!\big(X;\,E^q(\ptt)\big)\;\Longrightarrow\; E^{p+q}(X),
\end{equation}
converging to the filtration $F^pE^n(X)=\ker\!\big(E^n(X)\to E^n(X^{(p-1)})\big)$.
Here $E^q(\ptt)$ are the coefficient groups. When $E=H(-;A)$ is ordinary
cohomology this collapses at $E_2$; the content is entirely in \emph{generalized}
theories.

\subsection{Complex $K$--theory}

For $E=\KU$ (complex topological $K$--theory) Bott periodicity gives
\[
  \KU^q(\ptt)=\begin{cases}\Z,& q\ \text{even},\\ 0,& q\ \text{odd},\end{cases}
\]
with the Bott class $\beta_{\mathrm{Bott}}\in\KU^{-2}(\ptt)$ a unit. Hence
\begin{equation}\label{eq:E2K}
  E_2^{p,q}=\begin{cases}H^p(X;\Z),& q\ \text{even},\\ 0,& q\ \text{odd},\end{cases}
  \qquad\Longrightarrow\qquad \KU^{p+q}(X).
\end{equation}
The nonzero groups occupy the even rows; the whole picture is $2$--periodic in
$q$. Two immediate consequences:
\begin{itemize}
\item all \emph{even} differentials vanish, $d_{2k}=0$, because they map between
an even and an odd row (one of which is zero);
\item rationally the sequence degenerates, $E_2\otimes\Q\cong E_\infty\otimes\Q$,
so $\KU^0(X)\otimes\Q\cong\bigoplus_k H^{2k}(X;\Q)$ (the Chern character). All
higher differentials are therefore \emph{torsion} operations.
\end{itemize}

%======================================================================
\section{The first differential: $d_3=\widetilde{\Sq}^3$}
%======================================================================

\begin{thm}{Atiyah--Hirzebruch}{d3}
In the AHSS for $\KU$ the first potentially nonzero differential is
$d_3\colon E_3^{p,q}\to E_3^{p+3,q-2}$, and under \eqref{eq:E2K} it is the
\emph{integral Steenrod operation}
\[
  d_3=\widetilde{\Sq}^3\;=\;\beta\circ \Sq^2\circ \rho
  \;\colon\; H^{p}(X;\Z)\to H^{p+3}(X;\Z),
\]
where $\rho\colon H^*(-;\Z)\to H^*(-;\F_2)$ is reduction mod $2$, $\Sq^2$ is the
Steenrod square, and $\beta\colon H^*(-;\F_2)\to H^{*+1}(-;\Z)$ is the integral
Bockstein of $0\to\Z\xrightarrow{2}\Z\to\F_2\to 0$. In particular
$\widetilde{\Sq}^3$ takes values in $2$--torsion.
\end{thm}

\begin{titledBox}{Note II}{noteII}
Because $\rho\beta=\Sq^1$, the mod $2$ reduction of $\widetilde{\Sq}^3$ is the
genuine Steenrod square $\Sq^3=\Sq^1\Sq^2$. This is the fact we compute with in
Part~IV. Higher differentials $d_5,d_7,\dots$ exist and are secondary/tertiary
operations, but the counterexample only needs $d_3$.
\end{titledBox}

%======================================================================
\section{A second AHSS worked example: $\widetilde{\KU}{}^0(\RP^{2n})=\Z/2^n$}
%======================================================================

This example isolates the \emph{extension problem} in the AHSS (here every
differential vanishes, yet the answer is a nonsplit extension).

The integral cohomology of $X=\RP^{2n}$ is
\[
  H^{p}(\RP^{2n};\Z)=
  \begin{cases}\Z,& p=0,\\ \Z/2,& p=2,4,\dots,2n,\\ 0,& p\ \text{odd}.\end{cases}
\]
Plugging into \eqref{eq:E2K}, the $E_2$--page (row $q=0$; all other rows are a
$2$--periodic copy) has nonzero entries only in \emph{even} columns:
\[
\renewcommand{\arraystretch}{1.3}
\begin{array}{c|ccccccc}
 q=0 & \Z & 0 & \Z/2 & 0 & \Z/2 & \cdots & \Z/2\\ \hline
     & p=0 & 1 & 2 & 3 & 4 & \cdots & 2n
\end{array}
\]
Now run the checklist:
\begin{itemize}
\item \textbf{(D) Differentials.} $d_3\colon H^p\to H^{p+3}$ maps an even column
to an odd column, and $H^{\mathrm{odd}}(\RP^{2n})=0$. Hence $d_3=0$, and for the
same parity reason \emph{all} $d_r$ vanish. The sequence collapses:
$E_2=E_\infty$. (Note $\widetilde{\Sq}^3$ is identically $0$ here --- this space
is \emph{not} a counterexample; contrast Part~IV.)
\item \textbf{(E) Extensions.} The associated graded of
$\widetilde{\KU}{}^0(\RP^{2n})$ is $\bigoplus_{i=1}^{n}E_\infty^{2i,-2i}
=(\Z/2)^{\oplus n}$, of order $2^n$. But the filtration is nonsplit: writing
$\eta=\lambda_\C-1$ for the reduced class of the complexified tautological line
bundle, one has $\eta^{\,n+1}=0$ while $\eta$ has additive order $2^n$, so
\[
  \widetilde{\KU}{}^0(\RP^{2n})\;\cong\;\Z/2^n .
\]
The successive quotients $F^{2i}/F^{2i+2}=\Z/2$ assemble by \emph{nontrivial}
extensions $\Z/2^{i-1}\rightarrowtail \Z/2^{i}\twoheadrightarrow \Z/2$.
\end{itemize}
This is the same phenomenon as $a_1^2=2a_2$ in Part~I: $E_\infty$ knows the
graded pieces, not the group.

%======================================================================
\section{The obstruction: algebraic classes are permanent cycles}
%======================================================================

\subsection{Cycle classes lift to $K$--theory}

Let $X$ be a smooth complex projective variety and $Z\subset X$ a closed analytic
subset of complex codimension $k$. The coherent sheaf $\O_Z$ defines a class in
$K^0$ supported on $Z$; concretely a locally free resolution of $\O_Z$ gives a
class in $\KU^0(X)$ lying in the topological filtration step $F^{2k}$. Under the
Chern character / Grothendieck--Riemann--Roch this recovers the cycle class up to
a universal factor:
\begin{equation}\label{eq:GRR}
  c_k(\O_Z)=(-1)^{k-1}(k-1)!\,[Z]\in H^{2k}(X;\Z).
\end{equation}
The point is representability: $[Z]$ is (a rational multiple of) the image in
$\gr^{2k}\KU^0(X)=E_\infty^{2k,-2k}$ of an honest element of $\KU^0(X)$.

\begin{thm}{Atiyah--Hirzebruch criterion}{criterion}
For $Z\subset X$ closed analytic of codimension $k$, the cycle class
$[Z]\in H^{2k}(X;\Z)$ is a \emph{permanent cycle} in the AHSS for $\KU$: it lies
in the kernel of every differential $d_r$, $r\ge 2$. In particular
\[
  \widetilde{\Sq}^3\,[Z]=d_3[Z]=0 .
\]
More generally, all odd--degree Steenrod operations vanish on the classes of
algebraic cycles (in $H^*(X;\Z)$ and in $H^*(X;\Z/p)$).
\end{thm}

\begin{Proof}
\textbf{Symbol map.} The AHSS is the spectral sequence of the skeletal
filtration $F^p\KU^n(X)=\ker(\KU^n(X)\to\KU^n(X^{(p-1)}))$. Convergence gives a
surjection $\sigma_p\colon F^p\KU^n(X)\twoheadrightarrow F^p/F^{p+1}
=E_\infty^{p,n-p}$, and $E_\infty^{p,q}=Z_\infty^{p,q}/B_\infty^{p,q}$ is a
subquotient of $E_2^{p,q}$, where $Z_\infty=\bigcap_r\ker d_r$ (permanent cycles)
and $B_\infty=\bigcup_r\operatorname{im}d_r$ (boundaries), with
$B_\infty\subseteq Z_\infty$.

\smallskip
\textbf{Master lemma.} \emph{If $\alpha\in E_2^{p,q}$ equals $\sigma_p(\xi)\bmod
B_\infty$ for some $\xi\in F^p\KU^{p+q}(X)$, then $\alpha\in Z_\infty^{p,q}$.}
Indeed $\sigma_p(\xi)\in Z_\infty/B_\infty$, so the coset $\alpha+B_\infty$ lies
in the subgroup $Z_\infty/B_\infty$, whence $\alpha\in Z_\infty+B_\infty=Z_\infty$.
Conceptually: $d_r$ is the obstruction to extending a class from the
$(p{+}r{-}1)$--skeleton over the next one, and a globally defined $\xi$ extends
over every skeleton, so all $d_r$ vanish on its symbol.

\smallskip
\textbf{Geometric lift.} Let $i\colon Z\hookrightarrow X$ be the codimension--$k$
subvariety. Its complex normal data give a Gysin map
$i_!\colon\KU^0(Z)\to\KU^0(X)$; set $\xi=i_!(1)$. Since $i_!(1)$ is supported on
$Z$ and $\dim_{\mathbb R}Z=\dim_{\mathbb R}X-2k$, general position pushes
$X^{(2k-1)}$ off $Z$, so $i_!(1)|_{X^{(2k-1)}}=0$ and $\xi\in F^{2k}\KU^0(X)$.
By Riemann--Roch for $i_!$ (equivalently GRR), the symbol
$\sigma_{2k}(\xi)\in E_\infty^{2k,-2k}\subseteq_{\mathrm{s.q.}}H^{2k}(X;\Z)$ is the
Thom class of the normal bundle, i.e.\ the Poincar\'e dual $[Z]$. (Using $i_!(1)$
rather than $c_k(\mathcal O_Z)=(-1)^{k-1}(k-1)![Z]$ avoids the factorial, which
matters as $d_3$ is $2$--torsion valued; for $k=2$ the factor is $1$.)

\smallskip
\textbf{Conclusion.} The master lemma with $\alpha=[Z]$, $(p,q)=(2k,-2k)$ gives
$[Z]\in Z_\infty$, so $d_r[Z]=0$ for all $r$; in particular
$d_3[Z]=\widetilde{\Sq}^3[Z]=0$. The extension to all odd Steenrod operations
follows from the same argument in the mod--$p$ AHSS (Milnor primitives $Q_i$) and
Totaro's factorization through $MU$.
\end{Proof}

\begin{cor}{The contrapositive we exploit}{test}
If $\alpha\in H^{2k}(X;\Z)$ satisfies $\widetilde{\Sq}^3\alpha\neq 0$, then
$\alpha$ is \emph{not} in the image of the cycle class map
$\CH^k(X)\to H^{2k}(X;\Z)$; i.e.\ $\alpha$ is not algebraic. \qed
\end{cor}

%======================================================================
\section{The counterexample}
%======================================================================

\subsection{Statement}

\begin{defn}{Integral Hodge Classes}{intHodgeCls}
For $X$ smooth projective$/\C$, the group of \emph{integral Hodge classes} in
degree $2k$ is
\[
  \Hdg^{2k}(X,\Z)=\big\{\alpha\in H^{2k}(X;\Z):\ \alpha_\C\in H^{k,k}(X)\big\},
  \qquad \alpha_\C=\text{image in }H^{2k}(X;\C).
\]
The \emph{Integral Hodge Conjecture} (IHC) asserts that the cycle class map
$\CH^k(X)\to \Hdg^{2k}(X,\Z)$ is surjective.
\end{defn}

\begin{thm}{Atiyah--Hirzebruch, 1962}{AH}
The IHC is false for $k\ge 2$. Explicitly, there is a smooth complex projective
variety $X$ of dimension $7$ carrying a $2$--torsion class
$\alpha\in H^4(X;\Z)$ which is an integral Hodge class but is not the class of any
algebraic cycle.
\end{thm}

The proof has three independent ingredients: (i) torsion classes are automatically
Hodge; (ii) a construction of a projective $X$ realizing a chosen piece of group
cohomology; (iii) an explicit torsion class $\alpha$ with
$\widetilde{\Sq}^3\alpha\neq 0$, defeated by Corollary~\ref{co:test}.

\subsection{Torsion classes are Hodge}

\begin{lem}{torsion lies in Hodge}{torsionHodge}
Any torsion class $\alpha\in H^{2k}(X;\Z)_{\mathrm{tors}}$ lies in
$\Hdg^{2k}(X,\Z)$.
\end{lem}

\begin{proof}
The map $H^{2k}(X;\Z)\to H^{2k}(X;\C)$ kills torsion, since $H^{2k}(X;\C)$ is a
$\C$--vector space. Thus $\alpha_\C=0$, which lies in $H^{k,k}(X)$ trivially.
\end{proof}

So to break the IHC it suffices to produce a \emph{non‑algebraic torsion class} of
even degree.

\subsection{Realizing group cohomology projectively (Godeaux--Serre)}

\begin{thm}{Serre's construction; Atiyah--Hirzebruch Prop.~6.6}{GS}
Let $G$ be a finite group and $N\in\mathbb N$. Choose a faithful complex
representation $G\hookrightarrow GL(V)$; then $G$ acts freely on the complement of
a closed subvariety of high codimension in a large sum $V^{\oplus m}$, and
intersecting with generic hypersurfaces (Lefschetz) yields a smooth projective
variety $X$ with a free $G$--action such that the classifying map
$X/G\to BG$ induces isomorphisms
\[
  H^i(BG;\Z)\xrightarrow{\ \cong\ }H^i(X/G;\Z)\qquad(0\le i\le N).
\]
In fact $X/G\to BG\times BS^1$ is $r$--connected with $r=\dim_\C(X/G)$.
\end{thm}

Thus any torsion class in $H^{\le N}(BG;\Z)$ transports to a smooth projective
variety, together with all Steenrod operations, provided $\dim_\C X\ge N$. For our
class $\alpha\in H^4$ with $\widetilde{\Sq}^3\alpha\in H^7$, we need $N\ge 7$,
hence $\dim_\C X=7$ --- this is where the ``dimension $7$'' comes from.

\subsection{The explicit torsion class on $B(\Z/2)^3$}

Take $G=(\Z/2)^3$. With $\F_2$--coefficients,
\[
  H^*\!\big((\Z/2)^3;\F_2\big)=\F_2[x_1,x_2,x_3],\qquad |x_i|=1 .
\]
Let $w=x_1x_2x_3\in H^3$, and let $\alpha=\beta(w)\in H^4\big((\Z/2)^3;\Z\big)$ be
its integral Bockstein; $\alpha$ is $2$--torsion, and $\rho(\alpha)=\Sq^1 w$. We
show $\widetilde{\Sq}^3\alpha\neq 0$ by computing its mod $2$ reduction.

\smallskip\noindent\textbf{Reduction to $\Sq$'s.} Using $\rho\beta=\Sq^1$ and
Theorem~\ref{th:d3},
\[
  \rho\big(\widetilde{\Sq}^3\alpha\big)
  =\rho\,\beta\,\Sq^2\rho\,\beta(w)
  =\Sq^1\Sq^2\Sq^1(w).
\]

\smallskip\noindent\textbf{Adem relations.} From $\Sq^1\Sq^2=\Sq^3$ and
$\Sq^3\Sq^1=\Sq^2\Sq^2$ (equivalently $\Sq^2\Sq^2=\Sq^3\Sq^1$),
\[
  \Sq^1\Sq^2\Sq^1=\Sq^3\Sq^1=\Sq^2\Sq^2 .
\]
So it suffices to show $\Sq^2\Sq^2(w)\neq 0$.

\smallskip\noindent\textbf{The computation.} With the total square
$\Sq(x_i)=x_i+x_i^2$ and the Cartan formula, the degree‑$5$ part of $\Sq(w)$ is
\[
  \Sq^2(x_1x_2x_3)=x_1^2x_2^2x_3+x_1^2x_2x_3^2+x_1x_2^2x_3^2 .
\]
Applying $\Sq^2$ once more (using $\Sq(x_i^2)=x_i^2+x_i^4$) to each monomial and
adding mod $2$:
\[
\begin{aligned}
\Sq^2(x_1^2x_2^2x_3)&=x_1^4x_2^2x_3+x_1^2x_2^4x_3+x_1^2x_2^2x_3^2,\\
\Sq^2(x_1^2x_2x_3^2)&=x_1^4x_2x_3^2+x_1^2x_2^2x_3^2+x_1^2x_2x_3^4,\\
\Sq^2(x_1x_2^2x_3^2)&=x_1^2x_2^2x_3^2+x_1x_2^4x_3^2+x_1x_2^2x_3^4 .
\end{aligned}
\]
The three copies of $x_1^2x_2^2x_3^2$ add to a single copy (mod $2$), so
\[
  \boxed{\;\Sq^2\Sq^2(x_1x_2x_3)=x_1^2x_2^2x_3^2
  +\!\!\sum_{\text{$6$ perms of }(4,2,1)}\!\! x_1^{a}x_2^{b}x_3^{c}\;}
\]
\[
=x_1^2x_2^2x_3^2
+x_1^4x_2^2x_3+x_1^4x_2x_3^2+x_1^2x_2^4x_3+x_1x_2^4x_3^2+x_1^2x_2x_3^4+x_1x_2^2x_3^4 .
\]
These monomials are distinct and nonzero in the polynomial ring, so
$\Sq^2\Sq^2(w)\neq 0$. Consequently
\[
  \rho\big(\widetilde{\Sq}^3\alpha\big)=\Sq^1\Sq^2\Sq^1(w)=\Sq^2\Sq^2(w)\neq 0,
  \qquad\text{hence}\qquad \widetilde{\Sq}^3\alpha\neq 0 .
\]

\subsection{Conclusion}

\begin{proof}[Proof of Theorem~\ref{th:AH}]
By Theorem~\ref{th:GS} with $G=(\Z/2)^3$ and $N=7$, take a smooth projective
$7$--fold $X$ with $H^i(X;\Z)\cong H^i(B(\Z/2)^3;\Z)$ compatibly with Steenrod
operations for $i\le 7$. Let $\alpha\in H^4(X;\Z)$ be the image of
$\beta(x_1x_2x_3)$. Then:
\begin{itemize}
\item $\alpha$ is $2$--torsion, hence $\alpha\in\Hdg^4(X,\Z)$ by
Lemma~\ref{lm:torsionHodge};
\item $\widetilde{\Sq}^3\alpha\neq 0$ by the computation above (transported from
$B(\Z/2)^3$);
\item therefore $\alpha$ is not algebraic, by Corollary~\ref{co:test}.
\end{itemize}
So $\alpha$ is an integral Hodge class that is not a $\Z$--combination of
algebraic cycle classes: the IHC fails in codimension $2$.
\end{proof}

\begin{titledBox}{What ``core spectral sequence'' really means here}{noteIII}
The spectral sequence at the heart of the argument is the AHSS
\eqref{eq:E2K} for $\KU$. It enters not through a long page‑by‑page computation
but through one structural fact --- Theorem~\ref{th:criterion}: \emph{algebraic
classes are permanent cycles}. The obstruction is then localized in the single
differential $d_3=\widetilde{\Sq}^3$. In this sense the counterexample is a
one‑differential phenomenon, in contrast to the multi‑page bookkeeping of Part~I.
Totaro (1997) later refined the whole story by replacing $\KU$ with complex
cobordism $MU$, factoring the cycle map through $(MU^*(X)\otimes_{MU^*}\Z)^{2k}$
and producing sharper (higher‑differential) obstructions.
\end{titledBox}

%======================================================================
\begin{thebibliography}{9}
\bibitem{AH62} M.~F.~Atiyah and F.~Hirzebruch,
\emph{Analytic cycles on complex manifolds}, Topology \textbf{1} (1962), 25--45.
\bibitem{Tot97} B.~Totaro,
\emph{Torsion algebraic cycles and complex cobordism},
J.\ Amer.\ Math.\ Soc.\ \textbf{10} (1997), 467--493.
\bibitem{Voisin} C.~Voisin,
\emph{On the homotopy types of compact K\"ahler and complex projective manifolds}
and \emph{Some aspects of the Hodge conjecture}; see also her notes on torsion
cohomology classes and algebraic cycles.
\bibitem{McCleary} J.~McCleary,
\emph{A User's Guide to Spectral Sequences}, 2nd ed., Cambridge, 2001.
\bibitem{Hatcher} A.~Hatcher,
\emph{Spectral Sequences in Algebraic Topology} (online notes).
\end{thebibliography}

\end{document}
